\section{导读：一家车企为什么要从自动驾驶转向 Physical AI}

本节先说明这期的核心问题。刘先明接任小鹏自动驾驶中心负责人，外界关注的是换帅；但访谈真正值得整理的，是他把自动驾驶放进 Physical AI 的框架中重新理解：车不只是交通工具，也是一个真实世界智能体；自动驾驶不只是感知、规控和地图系统，而是数据、模型、云端工厂、车端部署和量产反馈共同构成的智能闭环。

这期最锋利的判断是“Language 是毒药”。它并不是说语言没有价值，而是说在自动驾驶/物理 AI 的训练链路中，如果把传感器信号先翻译成中间语言 token，再从语言解码轨迹，就会引入人工监督、压缩瓶颈和 data scaling 障碍。刘先明的路线是不断拆掉中间层：拆激光雷达依赖、拆规控规则、拆端到端中间结构，最后连 language bottleneck 也拆掉。

\begin{importantbox}{本期核心命题}
小鹏的 AI 转型不是“给汽车加大模型”，而是把自动驾驶重写成 Physical AI：用更大模型、更大数据、更少人工中间层和更强工程闭环，让车在真实世界中持续学习、部署和迭代。
\end{importantbox}

\begin{knowledgebox}{视觉策略说明}
本视频是固定访谈画面，没有 slides、白板或产品演示。正文只用封面做来源识别，正文图像全部为自制概念图，用来解释自动驾驶软件栈演化、拆掉 Language、云端模型工厂、主机厂数据闭环、组织扁平化和换帅使命。
\end{knowledgebox}

\subsection{本章小结}

EP120 的主线是“简化”：技术上拆掉中间层，组织上减少层级，战略上把自动驾驶放进 Physical AI。它不是单纯人事访谈，而是一份车企 AI 转型方法论。

